\chapter{Emergenz der Physik aus Q}

\section{Einleitung}

Die IWT ist die fundamentale Theorie der Information. Aus ihr emergiert die gesamte Physik – nicht willkürlich, sondern durch eine einzige, konsistente Ableitung.

In Kapitel 3 wurde gezeigt, dass aus dem Variationsprinzip der IWT eine Euler-Lagrange-Gleichung folgt, deren Lösung die Form \( I = \sqrt{\rho} \, e^{i\phi} \) hat. Die Separation in Real- und Imaginärteil liefert zwei Gleichungen:

\begin{enumerate}
\item Die \textbf{Realteil-Gleichung} – die Hamilton-Jacobi-Gleichung mit dem Bohm-Potential \( Q \).
\item Die \textbf{Imaginärteil-Gleichung} – die Kontinuitätsgleichung.
\end{enumerate}

Dieses Kapitel zeigt, dass aus derselben Euler-Lagrange-Gleichung – durch die Wahl des Real- oder Imaginärteils – die drei fundamentalen Theorien der Physik emergieren:

\begin{itemize}
\item die \textbf{De-Broglie-Bohm-Theorie (DBT)},
\item die \textbf{Weber-Gravitation (WG)},
\item die \textbf{Weber-Elektrodynamik (WED)}.
\end{itemize}

\section{Die gemeinsame Struktur}

Die Euler-Lagrange-Gleichung der IWT lautet:

\[
\frac{\partial \mathcal{L}_d}{\partial I} - \Delta_t^+ \left( \frac{\partial \mathcal{L}_d}{\partial (\Delta_t I)} \right) - \Delta_s \left( \frac{\partial \mathcal{L}_d}{\partial (\Delta_s I)} \right) = 0
\]

Setzt man \( I = \sqrt{\rho} \, e^{i\phi} \) ein und trennt Real- und Imaginärteil, so erhält man zwei unabhängige Gleichungen.

\subsection{Realteil: Die Hamilton-Jacobi-Gleichung}

Der Realteil der Euler-Lagrange-Gleichung ergibt:

\[
\frac{\partial S}{\partial t} + \frac{|\nabla S|^2}{2m} + V + Q = 0
\]

mit

\[
Q = -\frac{\hbar^2}{2m} \frac{\nabla^2 \sqrt{\rho}}{\sqrt{\rho}}
\]

Dabei ist \( S = \hbar \phi \). Dies ist die Hamilton-Jacobi-Gleichung der DBT.

\subsection{Imaginärteil: Die Kontinuitätsgleichung}

Der Imaginärteil der Euler-Lagrange-Gleichung ergibt:

\[
\frac{\partial \rho}{\partial t} + \nabla \cdot (\rho \vec{v}) = 0
\]

mit

\[
\vec{v} = \frac{\hbar}{m} \nabla \phi
\]

Dies ist die Kontinuitätsgleichung der DBT.

\section{Die Emergenz der DBT aus Q}

Die DBT ist die Quantentheorie.

Ihr Kern ist das Bohm-Potential:

\[
Q = -\frac{\hbar^2}{2m} \frac{\nabla^2 \sqrt{\rho}}{\sqrt{\rho}}
\]

Die DBT ist \( Q \). Sie emergiert direkt aus dem Realteil der IWT-Euler-Lagrange-Gleichung.

\section{Die Emergenz der WG aus Q}

Die Weber-Gravitation (WG) ist die Theorie der Gravitation.

Sie beschreibt eine direkte, geschwindigkeitsabhängige Wechselwirkung zwischen Massen.

Ihre Kraft lautet:

\[
\vec{F}_{\text{WG}} = -\frac{GMm}{r^2} \left[ 1 - \frac{\dot{r}^2}{c^2} + \frac{1}{2} \frac{r\ddot{r}}{c^2} \right] \hat{r}
\]

In der IWT emergiert die WG aus dem \textbf{Realteil} der Euler-Lagrange-Gleichung – also aus der \textbf{Amplitude} der Information.

Die Masse ist definiert als:

\[
m = \nabla |I|
\]

Die Gravitationskraft ist der negative Gradient des Potentials, das aus der Dichte der Masse folgt:

\[
\vec{F}_{\text{WG}} = -\nabla Q_{\text{WG}}
\]

mit

\[
Q_{\text{WG}} = -\frac{GM}{r}
\]

Die WG ist der Gradient von \( Q \), angewandt auf die Amplitude.

\section{Die Emergenz der WED aus Q}

Die Weber-Elektrodynamik (WED) ist die Theorie der Elektrodynamik.

Sie beschreibt eine direkte, geschwindigkeitsabhängige Wechselwirkung zwischen Ladungen.

Ihre Kraft lautet:

\[
\vec{F}_{\text{WED}} = \frac{q_1 q_2}{4\pi \varepsilon_0 r^2} \left[ 1 - \frac{\dot{r}^2}{c^2} + 2 \frac{r\ddot{r}}{c^2} \right] \hat{r}
\]

In der IWT emergiert die WED aus dem \textbf{Imaginärteil} der Euler-Lagrange-Gleichung – also aus der \textbf{Phase} der Information.

Die Ladung ist definiert als die eichinvariante Divergenz der Phase:

\[
q = \nabla^2 \phi
\]

Im diskreten Fall:

\[
q_k = \sum_{l \in \mathcal{N}(k)} (\phi_k - \phi_l)
\]

Die elektrische Kraft ist der negative Gradient des Potentials, das aus der Ladungsdichte folgt:

\[
\vec{F}_{\text{WED}} = -\nabla Q_{\text{WED}}
\]

mit

\[
Q_{\text{WED}} = \frac{q}{4\pi \varepsilon_0 r}
\]

Die WED ist der Gradient von \( Q \), angewandt auf die Phase.

\section{Die Quellen der drei Theorien}

\begin{itemize}
\item \textbf{DBT (Quantenmechanik):} Quelle ist die Wahrscheinlichkeitsdichte \( \rho = |I|^2 \). Das Bohm-Potential lautet \( Q = -\frac{\hbar^2}{2m} \frac{\nabla^2 \sqrt{\rho}}{\sqrt{\rho}} \).

\item \textbf{WG (Weber-Gravitation):} Quelle ist die Masse \( m = \nabla |I| \). Die Kraft ist \( \vec{F}_{\text{WG}} = -\nabla Q \) mit \( Q_{\text{WG}} = -\frac{GM}{r} \).

\item \textbf{WED (Weber-Elektrodynamik):} Quelle ist die Ladung \( q = \nabla^2 \phi \). Die Kraft ist \( \vec{F}_{\text{WED}} = -\nabla Q \) mit \( Q_{\text{WED}} = \frac{q}{4\pi \varepsilon_0 r} \).
\end{itemize}

\section{Die Einheit der Physik}

Die gesamte Physik ist \( Q \).

Die IWT ist die fundamentale Theorie. Aus ihr emergieren:

\begin{enumerate}
\item Die \textbf{DBT} aus dem Realteil der Euler-Lagrange-Gleichung – die Quantenmechanik.
\item Die \textbf{WG} aus dem Realteil – die Gravitation.
\item Die \textbf{WED} aus dem Imaginärteil – die Elektrodynamik.
\end{enumerate}

Die drei Theorien sind strukturell identisch:

\[
\begin{aligned}
\text{DBT} &= Q \\
\text{WG} &= -\nabla Q \quad (\text{angewandt auf die Amplitude}) \\
\text{WED} &= -\nabla Q \quad (\text{angewandt auf die Phase})
\end{aligned}
\]

\section{Die Konsequenz}

Die Physik ist vereinheitlicht.

Es gibt nur eine Theorie: die IWT.

Alles emergiert aus \( Q \).

\section{Zusammenfassung}

\begin{itemize}
\item Die DBT emergiert aus dem Realteil der IWT-Euler-Lagrange-Gleichung.
\item Die WG emergiert aus dem Realteil – als Gradient von \( Q \), angewandt auf die Masse (Amplitude).
\item Die WED emergiert aus dem Imaginärteil – als Gradient von \( Q \), angewandt auf die Ladung (Phase).
\item Die drei Theorien sind strukturell identisch.
\item Die Physik ist vereinheitlicht.
\end{itemize}